\documentclass[9pt]{extarticle}
\usepackage{core_attempt_forms}

\begin{document}

\CoreHeader{Counting}{C01}{A--E}{\FivePartScoring}
\begin{multicols}{2}
\begin{CoreProblem}{A}{Fundamental counting rules}
{\scriptsize\textit{Name and justify the model; exact fractions and unsimplified products are acceptable.}\par}
A cafe offers 5 soups and 4 salads. A lunch includes either one soup
    or one salad, but not both. How many lunches are possible? Name the counting
    rule.
\end{CoreProblem}
\begin{CoreProblem}{B}{Arrangements and selections}
{\scriptsize\textit{Name and justify the model; exact fractions and unsimplified products are acceptable.}\par}
A committee of 3 students is selected from 10 students. How many
    committees are possible? Give a binomial-coefficient expression and explain
    why order does not matter.
\end{CoreProblem}
\begin{CoreProblem}{C}{Constrained counting}
{\scriptsize\textit{Name and justify the model; exact fractions and unsimplified products are acceptable.}\par}
How many 5-card hands from a standard 52-card deck contain exactly
    two hearts? Give an unsimplified product of binomial coefficients.
\end{CoreProblem}
\columnbreak
\begin{CoreProblem}{D}{Probability models}
{\scriptsize\textit{Name and justify the model; exact fractions and unsimplified products are acceptable.}\par}
A fair coin is tossed and a fair six-sided die is rolled. Find the
    probability that the coin is heads or the die is even (or both). Show the
    favorable-outcome count or an equivalent probability calculation.
\end{CoreProblem}
\begin{CoreProblem}{E}{Modeling and interpretation}
{\scriptsize\textit{Name and justify the model; exact fractions and unsimplified products are acceptable.}\par}
Events $A$ and $B$ satisfy $P(A)=0.6$, $P(B)=0.5$, and
    $P(A\cap B)=0.3$. Determine whether $A$ and $B$ are independent, showing
    the numerical test, and then find $P(A\mid B)$.
\end{CoreProblem}
\end{multicols}
\CoreWorkArea
\FormBreak
\CoreHeader{Counting}{C02}{A--E}{\FivePartScoring}
\begin{multicols}{2}
\begin{CoreProblem}{A}{Fundamental counting rules}
{\scriptsize\textit{Name and justify the model; exact fractions and unsimplified products are acceptable.}\par}
A password contains two uppercase letters followed by three digits.
    Repetition is allowed. How many passwords are possible? Give an unsimplified
    product and name the counting rule.
\end{CoreProblem}
\begin{CoreProblem}{B}{Arrangements and selections}
{\scriptsize\textit{Name and justify the model; exact fractions and unsimplified products are acceptable.}\par}
A president and a vice president are selected from 10 students, with
    no student holding both offices. How many outcomes are possible? Give a
    permutation or product expression and explain why order matters.
\end{CoreProblem}
\begin{CoreProblem}{C}{Constrained counting}
{\scriptsize\textit{Name and justify the model; exact fractions and unsimplified products are acceptable.}\par}
How many bit strings of length 8 contain at least two $1$s? Use a
    complement and show the subtracted cases.
\end{CoreProblem}
\columnbreak
\begin{CoreProblem}{D}{Probability models}
{\scriptsize\textit{Name and justify the model; exact fractions and unsimplified products are acceptable.}\par}
Suppose $P(A)=0.55$, $P(B)=0.40$, and $P(A\cap B)=0.25$. Find
    $P(A\cup B)$ and $P(A^c\cap B)$, showing the probability identities used.
\end{CoreProblem}
\begin{CoreProblem}{E}{Modeling and interpretation}
{\scriptsize\textit{Name and justify the model; exact fractions and unsimplified products are acceptable.}\par}
A condition occurs in $1\%$ of a population. A test is positive for
    $95\%$ of people who have the condition and for $8\%$ of people who do not.
    Find $P(\text{positive})$ and then
    $P(\text{condition}\mid\text{positive})$. Give a one-sentence base-rate
    interpretation.
\end{CoreProblem}
\end{multicols}
\CoreWorkArea
\FormBreak
\CoreHeader{Counting}{C03}{A--E}{\FivePartScoring}
\begin{multicols}{2}
\begin{CoreProblem}{A}{Fundamental counting rules}
{\scriptsize\textit{Name and justify the model; exact fractions and unsimplified products are acceptable.}\par}
Let $A=\{1,2,3,4,5\}$ and $B=\{1,2,\ldots,8\}$. How many functions
    $f:A\to B$ exist? Identify the sequence of choices being counted.
\end{CoreProblem}
\begin{CoreProblem}{B}{Arrangements and selections}
{\scriptsize\textit{Name and justify the model; exact fractions and unsimplified products are acceptable.}\par}
How many distinct strings can be made by reordering the letters of
    \texttt{BANANA}? Give the appropriate multinomial expression.
\end{CoreProblem}
\begin{CoreProblem}{C}{Constrained counting}
{\scriptsize\textit{Name and justify the model; exact fractions and unsimplified products are acceptable.}\par}
A group of 5 students is selected from 12 computer-science majors and
    8 mathematics majors. How many groups contain exactly 3 computer-science
    majors? Give an unsimplified product of binomial coefficients.
\end{CoreProblem}
\columnbreak
\begin{CoreProblem}{D}{Probability models}
{\scriptsize\textit{Name and justify the model; exact fractions and unsimplified products are acceptable.}\par}
A uniformly chosen 8-bit string is selected. Find the probability
    that it contains exactly three $1$s. Give an exact fraction.
\end{CoreProblem}
\begin{CoreProblem}{E}{Modeling and interpretation}
{\scriptsize\textit{Name and justify the model; exact fractions and unsimplified products are acceptable.}\par}
Factory A supplies $40\%$ of a store's batteries and $3\%$ of its
    batteries are defective. Factory B supplies the rest and $1\%$ of its
    batteries are defective. A randomly selected battery is defective. Find the
    probability that it came from Factory A. Show total probability and Bayes'
    rule.
\end{CoreProblem}
\end{multicols}
\CoreWorkArea
\FormBreak
\CoreHeader{Counting}{C04}{A--E}{\FivePartScoring}
\begin{multicols}{2}
\begin{CoreProblem}{A}{Fundamental counting rules}
{\scriptsize\textit{Name and justify the model; exact fractions and unsimplified products are acceptable.}\par}
If $|A|=19$, $|B|=20$, and $|A\cap B|=8$, find $|A\cup B|$. Write
    the formula you used and name the principle.
\end{CoreProblem}
\begin{CoreProblem}{B}{Arrangements and selections}
{\scriptsize\textit{Name and justify the model; exact fractions and unsimplified products are acceptable.}\par}
How many bit strings of length 8 contain exactly three $1$s? Give a
    binomial-coefficient expression and identify what is being selected.
\end{CoreProblem}
\begin{CoreProblem}{C}{Constrained counting}
{\scriptsize\textit{Name and justify the model; exact fractions and unsimplified products are acceptable.}\par}
Let $|A|=5$ and $|B|=8$. How many injective functions $f:A\to B$
    exist? Give a permutation or product expression and explain the restriction.
\end{CoreProblem}
\columnbreak
\begin{CoreProblem}{D}{Probability models}
{\scriptsize\textit{Name and justify the model; exact fractions and unsimplified products are acceptable.}\par}
A bag contains 4 red and 6 blue tokens. Two tokens are drawn without
    replacement. Find the probability that the first is red and the second is
    blue. Show the conditional product.
\end{CoreProblem}
\begin{CoreProblem}{E}{Modeling and interpretation}
{\scriptsize\textit{Name and justify the model; exact fractions and unsimplified products are acceptable.}\par}
A game uses a fair six-sided die. It pays \$10 for a 6, \$2 for a 4
    or 5, and \$0 otherwise. A play costs \$2. Find the expected \emph{net}
    value of one play and interpret the sign.
\end{CoreProblem}
\end{multicols}
\CoreWorkArea
\FormBreak
\CoreHeader{Counting}{C05}{A--E}{\FivePartScoring}
\begin{multicols}{2}
\begin{CoreProblem}{A}{Fundamental counting rules}
{\scriptsize\textit{Name and justify the model; exact fractions and unsimplified products are acceptable.}\par}
A hash table maps 40 distinct keys to 20 possible hash values. What
    bucket occupancy is guaranteed? State the generalized pigeonhole calculation
    and explain why the data do not guarantee a larger occupancy.
\end{CoreProblem}
\begin{CoreProblem}{B}{Arrangements and selections}
{\scriptsize\textit{Name and justify the model; exact fractions and unsimplified products are acceptable.}\par}
Find the coefficient of $x^3$ in $(1+x)^8$. State the binomial-theorem
    term that determines your answer.
\end{CoreProblem}
\begin{CoreProblem}{C}{Constrained counting}
{\scriptsize\textit{Name and justify the model; exact fractions and unsimplified products are acceptable.}\par}
A six-character code uses uppercase letters without repetition and
    must begin with a vowel. How many codes are possible? Treat
    $A,E,I,O,U$ as the five vowels and give an unsimplified product.
\end{CoreProblem}
\columnbreak
\begin{CoreProblem}{D}{Probability models}
{\scriptsize\textit{Name and justify the model; exact fractions and unsimplified products are acceptable.}\par}
In a class, 24 students use Python, 18 use Java, and 10 use both.
    One student is chosen uniformly. Find
    $P(\text{Python}\mid\text{Java})$ and
    $P(\text{Java}\mid\text{Python})$.
\end{CoreProblem}
\begin{CoreProblem}{E}{Modeling and interpretation}
{\scriptsize\textit{Name and justify the model; exact fractions and unsimplified products are acceptable.}\par}
The table gives successes/participants for two treatments.
    \[
    \begin{array}{c|cc}
      &\text{Treatment A}&\text{Treatment B}\\\hline
      \text{Group 1}&30/100&400/1000\\
      \text{Group 2}&800/1000&90/100
    \end{array}
    \]
    Identify the better treatment within each group and overall. Then explain in
    one sentence how the unequal treatment denominators produce the reversal.
\end{CoreProblem}
\end{multicols}
\CoreWorkArea
\FormBreak
\CoreHeader{Counting}{C06}{A--E}{\FivePartScoring}
\begin{multicols}{2}
\begin{CoreProblem}{A}{Fundamental counting rules}
{\scriptsize\textit{Name and justify the model; exact fractions and unsimplified products are acceptable.}\par}
A cafe offers 5 soups and 4 salads. A lunch includes either one soup
    or one salad, but not both. How many lunches are possible? Name the counting
    rule.
\end{CoreProblem}
\begin{CoreProblem}{B}{Arrangements and selections}
{\scriptsize\textit{Name and justify the model; exact fractions and unsimplified products are acceptable.}\par}
A president and a vice president are selected from 10 students, with
    no student holding both offices. How many outcomes are possible? Give a
    permutation or product expression and explain why order matters.
\end{CoreProblem}
\begin{CoreProblem}{C}{Constrained counting}
{\scriptsize\textit{Name and justify the model; exact fractions and unsimplified products are acceptable.}\par}
A group of 5 students is selected from 12 computer-science majors and
    8 mathematics majors. How many groups contain exactly 3 computer-science
    majors? Give an unsimplified product of binomial coefficients.
\end{CoreProblem}
\columnbreak
\begin{CoreProblem}{D}{Probability models}
{\scriptsize\textit{Name and justify the model; exact fractions and unsimplified products are acceptable.}\par}
A bag contains 4 red and 6 blue tokens. Two tokens are drawn without
    replacement. Find the probability that the first is red and the second is
    blue. Show the conditional product.
\end{CoreProblem}
\begin{CoreProblem}{E}{Modeling and interpretation}
{\scriptsize\textit{Name and justify the model; exact fractions and unsimplified products are acceptable.}\par}
The table gives successes/participants for two treatments.
    \[
    \begin{array}{c|cc}
      &\text{Treatment A}&\text{Treatment B}\\\hline
      \text{Group 1}&30/100&400/1000\\
      \text{Group 2}&800/1000&90/100
    \end{array}
    \]
    Identify the better treatment within each group and overall. Then explain in
    one sentence how the unequal treatment denominators produce the reversal.
\end{CoreProblem}
\end{multicols}
\CoreWorkArea
\FormBreak
\CoreHeader{Counting}{C07}{A--E}{\FivePartScoring}
\begin{multicols}{2}
\begin{CoreProblem}{A}{Fundamental counting rules}
{\scriptsize\textit{Name and justify the model; exact fractions and unsimplified products are acceptable.}\par}
A password contains two uppercase letters followed by three digits.
    Repetition is allowed. How many passwords are possible? Give an unsimplified
    product and name the counting rule.
\end{CoreProblem}
\begin{CoreProblem}{B}{Arrangements and selections}
{\scriptsize\textit{Name and justify the model; exact fractions and unsimplified products are acceptable.}\par}
How many distinct strings can be made by reordering the letters of
    \texttt{BANANA}? Give the appropriate multinomial expression.
\end{CoreProblem}
\begin{CoreProblem}{C}{Constrained counting}
{\scriptsize\textit{Name and justify the model; exact fractions and unsimplified products are acceptable.}\par}
Let $|A|=5$ and $|B|=8$. How many injective functions $f:A\to B$
    exist? Give a permutation or product expression and explain the restriction.
\end{CoreProblem}
\columnbreak
\begin{CoreProblem}{D}{Probability models}
{\scriptsize\textit{Name and justify the model; exact fractions and unsimplified products are acceptable.}\par}
In a class, 24 students use Python, 18 use Java, and 10 use both.
    One student is chosen uniformly. Find
    $P(\text{Python}\mid\text{Java})$ and
    $P(\text{Java}\mid\text{Python})$.
\end{CoreProblem}
\begin{CoreProblem}{E}{Modeling and interpretation}
{\scriptsize\textit{Name and justify the model; exact fractions and unsimplified products are acceptable.}\par}
Events $A$ and $B$ satisfy $P(A)=0.6$, $P(B)=0.5$, and
    $P(A\cap B)=0.3$. Determine whether $A$ and $B$ are independent, showing
    the numerical test, and then find $P(A\mid B)$.
\end{CoreProblem}
\end{multicols}
\CoreWorkArea
\FormBreak
\CoreHeader{Counting}{C08}{A--E}{\FivePartScoring}
\begin{multicols}{2}
\begin{CoreProblem}{A}{Fundamental counting rules}
{\scriptsize\textit{Name and justify the model; exact fractions and unsimplified products are acceptable.}\par}
Let $A=\{1,2,3,4,5\}$ and $B=\{1,2,\ldots,8\}$. How many functions
    $f:A\to B$ exist? Identify the sequence of choices being counted.
\end{CoreProblem}
\begin{CoreProblem}{B}{Arrangements and selections}
{\scriptsize\textit{Name and justify the model; exact fractions and unsimplified products are acceptable.}\par}
How many bit strings of length 8 contain exactly three $1$s? Give a
    binomial-coefficient expression and identify what is being selected.
\end{CoreProblem}
\begin{CoreProblem}{C}{Constrained counting}
{\scriptsize\textit{Name and justify the model; exact fractions and unsimplified products are acceptable.}\par}
A six-character code uses uppercase letters without repetition and
    must begin with a vowel. How many codes are possible? Treat
    $A,E,I,O,U$ as the five vowels and give an unsimplified product.
\end{CoreProblem}
\columnbreak
\begin{CoreProblem}{D}{Probability models}
{\scriptsize\textit{Name and justify the model; exact fractions and unsimplified products are acceptable.}\par}
A fair coin is tossed and a fair six-sided die is rolled. Find the
    probability that the coin is heads or the die is even (or both). Show the
    favorable-outcome count or an equivalent probability calculation.
\end{CoreProblem}
\begin{CoreProblem}{E}{Modeling and interpretation}
{\scriptsize\textit{Name and justify the model; exact fractions and unsimplified products are acceptable.}\par}
A condition occurs in $1\%$ of a population. A test is positive for
    $95\%$ of people who have the condition and for $8\%$ of people who do not.
    Find $P(\text{positive})$ and then
    $P(\text{condition}\mid\text{positive})$. Give a one-sentence base-rate
    interpretation.
\end{CoreProblem}
\end{multicols}
\CoreWorkArea
\FormBreak
\CoreHeader{Counting}{C09}{A--E}{\FivePartScoring}
\begin{multicols}{2}
\begin{CoreProblem}{A}{Fundamental counting rules}
{\scriptsize\textit{Name and justify the model; exact fractions and unsimplified products are acceptable.}\par}
If $|A|=19$, $|B|=20$, and $|A\cap B|=8$, find $|A\cup B|$. Write
    the formula you used and name the principle.
\end{CoreProblem}
\begin{CoreProblem}{B}{Arrangements and selections}
{\scriptsize\textit{Name and justify the model; exact fractions and unsimplified products are acceptable.}\par}
Find the coefficient of $x^3$ in $(1+x)^8$. State the binomial-theorem
    term that determines your answer.
\end{CoreProblem}
\begin{CoreProblem}{C}{Constrained counting}
{\scriptsize\textit{Name and justify the model; exact fractions and unsimplified products are acceptable.}\par}
How many 5-card hands from a standard 52-card deck contain exactly
    two hearts? Give an unsimplified product of binomial coefficients.
\end{CoreProblem}
\columnbreak
\begin{CoreProblem}{D}{Probability models}
{\scriptsize\textit{Name and justify the model; exact fractions and unsimplified products are acceptable.}\par}
Suppose $P(A)=0.55$, $P(B)=0.40$, and $P(A\cap B)=0.25$. Find
    $P(A\cup B)$ and $P(A^c\cap B)$, showing the probability identities used.
\end{CoreProblem}
\begin{CoreProblem}{E}{Modeling and interpretation}
{\scriptsize\textit{Name and justify the model; exact fractions and unsimplified products are acceptable.}\par}
Factory A supplies $40\%$ of a store's batteries and $3\%$ of its
    batteries are defective. Factory B supplies the rest and $1\%$ of its
    batteries are defective. A randomly selected battery is defective. Find the
    probability that it came from Factory A. Show total probability and Bayes'
    rule.
\end{CoreProblem}
\end{multicols}
\CoreWorkArea
\FormBreak
\CoreHeader{Counting}{C10}{A--E}{\FivePartScoring}
\begin{multicols}{2}
\begin{CoreProblem}{A}{Fundamental counting rules}
{\scriptsize\textit{Name and justify the model; exact fractions and unsimplified products are acceptable.}\par}
A hash table maps 40 distinct keys to 20 possible hash values. What
    bucket occupancy is guaranteed? State the generalized pigeonhole calculation
    and explain why the data do not guarantee a larger occupancy.
\end{CoreProblem}
\begin{CoreProblem}{B}{Arrangements and selections}
{\scriptsize\textit{Name and justify the model; exact fractions and unsimplified products are acceptable.}\par}
A committee of 3 students is selected from 10 students. How many
    committees are possible? Give a binomial-coefficient expression and explain
    why order does not matter.
\end{CoreProblem}
\begin{CoreProblem}{C}{Constrained counting}
{\scriptsize\textit{Name and justify the model; exact fractions and unsimplified products are acceptable.}\par}
How many bit strings of length 8 contain at least two $1$s? Use a
    complement and show the subtracted cases.
\end{CoreProblem}
\columnbreak
\begin{CoreProblem}{D}{Probability models}
{\scriptsize\textit{Name and justify the model; exact fractions and unsimplified products are acceptable.}\par}
A uniformly chosen 8-bit string is selected. Find the probability
    that it contains exactly three $1$s. Give an exact fraction.
\end{CoreProblem}
\begin{CoreProblem}{E}{Modeling and interpretation}
{\scriptsize\textit{Name and justify the model; exact fractions and unsimplified products are acceptable.}\par}
A game uses a fair six-sided die. It pays \$10 for a 6, \$2 for a 4
    or 5, and \$0 otherwise. A play costs \$2. Find the expected \emph{net}
    value of one play and interpret the sign.
\end{CoreProblem}
\end{multicols}
\CoreWorkArea
\FormBreak
\CoreHeader{Counting}{C11}{A--E}{\FivePartScoring}
\begin{multicols}{2}
\begin{CoreProblem}{A}{Fundamental counting rules}
{\scriptsize\textit{Name and justify the model; exact fractions and unsimplified products are acceptable.}\par}
A cafe offers 5 soups and 4 salads. A lunch includes either one soup
    or one salad, but not both. How many lunches are possible? Name the counting
    rule.
\end{CoreProblem}
\begin{CoreProblem}{B}{Arrangements and selections}
{\scriptsize\textit{Name and justify the model; exact fractions and unsimplified products are acceptable.}\par}
How many distinct strings can be made by reordering the letters of
    \texttt{BANANA}? Give the appropriate multinomial expression.
\end{CoreProblem}
\begin{CoreProblem}{C}{Constrained counting}
{\scriptsize\textit{Name and justify the model; exact fractions and unsimplified products are acceptable.}\par}
A six-character code uses uppercase letters without repetition and
    must begin with a vowel. How many codes are possible? Treat
    $A,E,I,O,U$ as the five vowels and give an unsimplified product.
\end{CoreProblem}
\columnbreak
\begin{CoreProblem}{D}{Probability models}
{\scriptsize\textit{Name and justify the model; exact fractions and unsimplified products are acceptable.}\par}
Suppose $P(A)=0.55$, $P(B)=0.40$, and $P(A\cap B)=0.25$. Find
    $P(A\cup B)$ and $P(A^c\cap B)$, showing the probability identities used.
\end{CoreProblem}
\begin{CoreProblem}{E}{Modeling and interpretation}
{\scriptsize\textit{Name and justify the model; exact fractions and unsimplified products are acceptable.}\par}
A game uses a fair six-sided die. It pays \$10 for a 6, \$2 for a 4
    or 5, and \$0 otherwise. A play costs \$2. Find the expected \emph{net}
    value of one play and interpret the sign.
\end{CoreProblem}
\end{multicols}
\CoreWorkArea
\FormBreak
\CoreHeader{Counting}{C12}{A--E}{\FivePartScoring}
\begin{multicols}{2}
\begin{CoreProblem}{A}{Fundamental counting rules}
{\scriptsize\textit{Name and justify the model; exact fractions and unsimplified products are acceptable.}\par}
A password contains two uppercase letters followed by three digits.
    Repetition is allowed. How many passwords are possible? Give an unsimplified
    product and name the counting rule.
\end{CoreProblem}
\begin{CoreProblem}{B}{Arrangements and selections}
{\scriptsize\textit{Name and justify the model; exact fractions and unsimplified products are acceptable.}\par}
How many bit strings of length 8 contain exactly three $1$s? Give a
    binomial-coefficient expression and identify what is being selected.
\end{CoreProblem}
\begin{CoreProblem}{C}{Constrained counting}
{\scriptsize\textit{Name and justify the model; exact fractions and unsimplified products are acceptable.}\par}
How many 5-card hands from a standard 52-card deck contain exactly
    two hearts? Give an unsimplified product of binomial coefficients.
\end{CoreProblem}
\columnbreak
\begin{CoreProblem}{D}{Probability models}
{\scriptsize\textit{Name and justify the model; exact fractions and unsimplified products are acceptable.}\par}
A uniformly chosen 8-bit string is selected. Find the probability
    that it contains exactly three $1$s. Give an exact fraction.
\end{CoreProblem}
\begin{CoreProblem}{E}{Modeling and interpretation}
{\scriptsize\textit{Name and justify the model; exact fractions and unsimplified products are acceptable.}\par}
The table gives successes/participants for two treatments.
    \[
    \begin{array}{c|cc}
      &\text{Treatment A}&\text{Treatment B}\\\hline
      \text{Group 1}&30/100&400/1000\\
      \text{Group 2}&800/1000&90/100
    \end{array}
    \]
    Identify the better treatment within each group and overall. Then explain in
    one sentence how the unequal treatment denominators produce the reversal.
\end{CoreProblem}
\end{multicols}
\CoreWorkArea
\FormBreak
\CoreHeader{Counting}{C13}{A--E}{\FivePartScoring}
\begin{multicols}{2}
\begin{CoreProblem}{A}{Fundamental counting rules}
{\scriptsize\textit{Name and justify the model; exact fractions and unsimplified products are acceptable.}\par}
Let $A=\{1,2,3,4,5\}$ and $B=\{1,2,\ldots,8\}$. How many functions
    $f:A\to B$ exist? Identify the sequence of choices being counted.
\end{CoreProblem}
\begin{CoreProblem}{B}{Arrangements and selections}
{\scriptsize\textit{Name and justify the model; exact fractions and unsimplified products are acceptable.}\par}
Find the coefficient of $x^3$ in $(1+x)^8$. State the binomial-theorem
    term that determines your answer.
\end{CoreProblem}
\begin{CoreProblem}{C}{Constrained counting}
{\scriptsize\textit{Name and justify the model; exact fractions and unsimplified products are acceptable.}\par}
How many bit strings of length 8 contain at least two $1$s? Use a
    complement and show the subtracted cases.
\end{CoreProblem}
\columnbreak
\begin{CoreProblem}{D}{Probability models}
{\scriptsize\textit{Name and justify the model; exact fractions and unsimplified products are acceptable.}\par}
A bag contains 4 red and 6 blue tokens. Two tokens are drawn without
    replacement. Find the probability that the first is red and the second is
    blue. Show the conditional product.
\end{CoreProblem}
\begin{CoreProblem}{E}{Modeling and interpretation}
{\scriptsize\textit{Name and justify the model; exact fractions and unsimplified products are acceptable.}\par}
Events $A$ and $B$ satisfy $P(A)=0.6$, $P(B)=0.5$, and
    $P(A\cap B)=0.3$. Determine whether $A$ and $B$ are independent, showing
    the numerical test, and then find $P(A\mid B)$.
\end{CoreProblem}
\end{multicols}
\CoreWorkArea
\FormBreak
\CoreHeader{Counting}{C14}{A--E}{\FivePartScoring}
\begin{multicols}{2}
\begin{CoreProblem}{A}{Fundamental counting rules}
{\scriptsize\textit{Name and justify the model; exact fractions and unsimplified products are acceptable.}\par}
If $|A|=19$, $|B|=20$, and $|A\cap B|=8$, find $|A\cup B|$. Write
    the formula you used and name the principle.
\end{CoreProblem}
\begin{CoreProblem}{B}{Arrangements and selections}
{\scriptsize\textit{Name and justify the model; exact fractions and unsimplified products are acceptable.}\par}
A committee of 3 students is selected from 10 students. How many
    committees are possible? Give a binomial-coefficient expression and explain
    why order does not matter.
\end{CoreProblem}
\begin{CoreProblem}{C}{Constrained counting}
{\scriptsize\textit{Name and justify the model; exact fractions and unsimplified products are acceptable.}\par}
A group of 5 students is selected from 12 computer-science majors and
    8 mathematics majors. How many groups contain exactly 3 computer-science
    majors? Give an unsimplified product of binomial coefficients.
\end{CoreProblem}
\columnbreak
\begin{CoreProblem}{D}{Probability models}
{\scriptsize\textit{Name and justify the model; exact fractions and unsimplified products are acceptable.}\par}
In a class, 24 students use Python, 18 use Java, and 10 use both.
    One student is chosen uniformly. Find
    $P(\text{Python}\mid\text{Java})$ and
    $P(\text{Java}\mid\text{Python})$.
\end{CoreProblem}
\begin{CoreProblem}{E}{Modeling and interpretation}
{\scriptsize\textit{Name and justify the model; exact fractions and unsimplified products are acceptable.}\par}
A condition occurs in $1\%$ of a population. A test is positive for
    $95\%$ of people who have the condition and for $8\%$ of people who do not.
    Find $P(\text{positive})$ and then
    $P(\text{condition}\mid\text{positive})$. Give a one-sentence base-rate
    interpretation.
\end{CoreProblem}
\end{multicols}
\CoreWorkArea
\FormBreak
\CoreHeader{Counting}{C15}{A--E}{\FivePartScoring}
\begin{multicols}{2}
\begin{CoreProblem}{A}{Fundamental counting rules}
{\scriptsize\textit{Name and justify the model; exact fractions and unsimplified products are acceptable.}\par}
A hash table maps 40 distinct keys to 20 possible hash values. What
    bucket occupancy is guaranteed? State the generalized pigeonhole calculation
    and explain why the data do not guarantee a larger occupancy.
\end{CoreProblem}
\begin{CoreProblem}{B}{Arrangements and selections}
{\scriptsize\textit{Name and justify the model; exact fractions and unsimplified products are acceptable.}\par}
A president and a vice president are selected from 10 students, with
    no student holding both offices. How many outcomes are possible? Give a
    permutation or product expression and explain why order matters.
\end{CoreProblem}
\begin{CoreProblem}{C}{Constrained counting}
{\scriptsize\textit{Name and justify the model; exact fractions and unsimplified products are acceptable.}\par}
Let $|A|=5$ and $|B|=8$. How many injective functions $f:A\to B$
    exist? Give a permutation or product expression and explain the restriction.
\end{CoreProblem}
\columnbreak
\begin{CoreProblem}{D}{Probability models}
{\scriptsize\textit{Name and justify the model; exact fractions and unsimplified products are acceptable.}\par}
A fair coin is tossed and a fair six-sided die is rolled. Find the
    probability that the coin is heads or the die is even (or both). Show the
    favorable-outcome count or an equivalent probability calculation.
\end{CoreProblem}
\begin{CoreProblem}{E}{Modeling and interpretation}
{\scriptsize\textit{Name and justify the model; exact fractions and unsimplified products are acceptable.}\par}
Factory A supplies $40\%$ of a store's batteries and $3\%$ of its
    batteries are defective. Factory B supplies the rest and $1\%$ of its
    batteries are defective. A randomly selected battery is defective. Find the
    probability that it came from Factory A. Show total probability and Bayes'
    rule.
\end{CoreProblem}
\end{multicols}
\CoreWorkArea
\FormBreak
\CoreHeader{Counting}{C16}{A--E}{\FivePartScoring}
\begin{multicols}{2}
\begin{CoreProblem}{A}{Fundamental counting rules}
{\scriptsize\textit{Name and justify the model; exact fractions and unsimplified products are acceptable.}\par}
A cafe offers 5 soups and 4 salads. A lunch includes either one soup
    or one salad, but not both. How many lunches are possible? Name the counting
    rule.
\end{CoreProblem}
\begin{CoreProblem}{B}{Arrangements and selections}
{\scriptsize\textit{Name and justify the model; exact fractions and unsimplified products are acceptable.}\par}
How many bit strings of length 8 contain exactly three $1$s? Give a
    binomial-coefficient expression and identify what is being selected.
\end{CoreProblem}
\begin{CoreProblem}{C}{Constrained counting}
{\scriptsize\textit{Name and justify the model; exact fractions and unsimplified products are acceptable.}\par}
How many bit strings of length 8 contain at least two $1$s? Use a
    complement and show the subtracted cases.
\end{CoreProblem}
\columnbreak
\begin{CoreProblem}{D}{Probability models}
{\scriptsize\textit{Name and justify the model; exact fractions and unsimplified products are acceptable.}\par}
In a class, 24 students use Python, 18 use Java, and 10 use both.
    One student is chosen uniformly. Find
    $P(\text{Python}\mid\text{Java})$ and
    $P(\text{Java}\mid\text{Python})$.
\end{CoreProblem}
\begin{CoreProblem}{E}{Modeling and interpretation}
{\scriptsize\textit{Name and justify the model; exact fractions and unsimplified products are acceptable.}\par}
Factory A supplies $40\%$ of a store's batteries and $3\%$ of its
    batteries are defective. Factory B supplies the rest and $1\%$ of its
    batteries are defective. A randomly selected battery is defective. Find the
    probability that it came from Factory A. Show total probability and Bayes'
    rule.
\end{CoreProblem}
\end{multicols}
\CoreWorkArea
\FormBreak
\CoreHeader{Counting}{C17}{A--E}{\FivePartScoring}
\begin{multicols}{2}
\begin{CoreProblem}{A}{Fundamental counting rules}
{\scriptsize\textit{Name and justify the model; exact fractions and unsimplified products are acceptable.}\par}
A password contains two uppercase letters followed by three digits.
    Repetition is allowed. How many passwords are possible? Give an unsimplified
    product and name the counting rule.
\end{CoreProblem}
\begin{CoreProblem}{B}{Arrangements and selections}
{\scriptsize\textit{Name and justify the model; exact fractions and unsimplified products are acceptable.}\par}
Find the coefficient of $x^3$ in $(1+x)^8$. State the binomial-theorem
    term that determines your answer.
\end{CoreProblem}
\begin{CoreProblem}{C}{Constrained counting}
{\scriptsize\textit{Name and justify the model; exact fractions and unsimplified products are acceptable.}\par}
A group of 5 students is selected from 12 computer-science majors and
    8 mathematics majors. How many groups contain exactly 3 computer-science
    majors? Give an unsimplified product of binomial coefficients.
\end{CoreProblem}
\columnbreak
\begin{CoreProblem}{D}{Probability models}
{\scriptsize\textit{Name and justify the model; exact fractions and unsimplified products are acceptable.}\par}
A fair coin is tossed and a fair six-sided die is rolled. Find the
    probability that the coin is heads or the die is even (or both). Show the
    favorable-outcome count or an equivalent probability calculation.
\end{CoreProblem}
\begin{CoreProblem}{E}{Modeling and interpretation}
{\scriptsize\textit{Name and justify the model; exact fractions and unsimplified products are acceptable.}\par}
A game uses a fair six-sided die. It pays \$10 for a 6, \$2 for a 4
    or 5, and \$0 otherwise. A play costs \$2. Find the expected \emph{net}
    value of one play and interpret the sign.
\end{CoreProblem}
\end{multicols}
\CoreWorkArea
\FormBreak
\CoreHeader{Counting}{C18}{A--E}{\FivePartScoring}
\begin{multicols}{2}
\begin{CoreProblem}{A}{Fundamental counting rules}
{\scriptsize\textit{Name and justify the model; exact fractions and unsimplified products are acceptable.}\par}
Let $A=\{1,2,3,4,5\}$ and $B=\{1,2,\ldots,8\}$. How many functions
    $f:A\to B$ exist? Identify the sequence of choices being counted.
\end{CoreProblem}
\begin{CoreProblem}{B}{Arrangements and selections}
{\scriptsize\textit{Name and justify the model; exact fractions and unsimplified products are acceptable.}\par}
A committee of 3 students is selected from 10 students. How many
    committees are possible? Give a binomial-coefficient expression and explain
    why order does not matter.
\end{CoreProblem}
\begin{CoreProblem}{C}{Constrained counting}
{\scriptsize\textit{Name and justify the model; exact fractions and unsimplified products are acceptable.}\par}
Let $|A|=5$ and $|B|=8$. How many injective functions $f:A\to B$
    exist? Give a permutation or product expression and explain the restriction.
\end{CoreProblem}
\columnbreak
\begin{CoreProblem}{D}{Probability models}
{\scriptsize\textit{Name and justify the model; exact fractions and unsimplified products are acceptable.}\par}
Suppose $P(A)=0.55$, $P(B)=0.40$, and $P(A\cap B)=0.25$. Find
    $P(A\cup B)$ and $P(A^c\cap B)$, showing the probability identities used.
\end{CoreProblem}
\begin{CoreProblem}{E}{Modeling and interpretation}
{\scriptsize\textit{Name and justify the model; exact fractions and unsimplified products are acceptable.}\par}
The table gives successes/participants for two treatments.
    \[
    \begin{array}{c|cc}
      &\text{Treatment A}&\text{Treatment B}\\\hline
      \text{Group 1}&30/100&400/1000\\
      \text{Group 2}&800/1000&90/100
    \end{array}
    \]
    Identify the better treatment within each group and overall. Then explain in
    one sentence how the unequal treatment denominators produce the reversal.
\end{CoreProblem}
\end{multicols}
\CoreWorkArea
\FormBreak
\CoreHeader{Counting}{C19}{A--E}{\FivePartScoring}
\begin{multicols}{2}
\begin{CoreProblem}{A}{Fundamental counting rules}
{\scriptsize\textit{Name and justify the model; exact fractions and unsimplified products are acceptable.}\par}
If $|A|=19$, $|B|=20$, and $|A\cap B|=8$, find $|A\cup B|$. Write
    the formula you used and name the principle.
\end{CoreProblem}
\begin{CoreProblem}{B}{Arrangements and selections}
{\scriptsize\textit{Name and justify the model; exact fractions and unsimplified products are acceptable.}\par}
A president and a vice president are selected from 10 students, with
    no student holding both offices. How many outcomes are possible? Give a
    permutation or product expression and explain why order matters.
\end{CoreProblem}
\begin{CoreProblem}{C}{Constrained counting}
{\scriptsize\textit{Name and justify the model; exact fractions and unsimplified products are acceptable.}\par}
A six-character code uses uppercase letters without repetition and
    must begin with a vowel. How many codes are possible? Treat
    $A,E,I,O,U$ as the five vowels and give an unsimplified product.
\end{CoreProblem}
\columnbreak
\begin{CoreProblem}{D}{Probability models}
{\scriptsize\textit{Name and justify the model; exact fractions and unsimplified products are acceptable.}\par}
A uniformly chosen 8-bit string is selected. Find the probability
    that it contains exactly three $1$s. Give an exact fraction.
\end{CoreProblem}
\begin{CoreProblem}{E}{Modeling and interpretation}
{\scriptsize\textit{Name and justify the model; exact fractions and unsimplified products are acceptable.}\par}
Events $A$ and $B$ satisfy $P(A)=0.6$, $P(B)=0.5$, and
    $P(A\cap B)=0.3$. Determine whether $A$ and $B$ are independent, showing
    the numerical test, and then find $P(A\mid B)$.
\end{CoreProblem}
\end{multicols}
\CoreWorkArea
\FormBreak
\CoreHeader{Counting}{C20}{A--E}{\FivePartScoring}
\begin{multicols}{2}
\begin{CoreProblem}{A}{Fundamental counting rules}
{\scriptsize\textit{Name and justify the model; exact fractions and unsimplified products are acceptable.}\par}
A hash table maps 40 distinct keys to 20 possible hash values. What
    bucket occupancy is guaranteed? State the generalized pigeonhole calculation
    and explain why the data do not guarantee a larger occupancy.
\end{CoreProblem}
\begin{CoreProblem}{B}{Arrangements and selections}
{\scriptsize\textit{Name and justify the model; exact fractions and unsimplified products are acceptable.}\par}
How many distinct strings can be made by reordering the letters of
    \texttt{BANANA}? Give the appropriate multinomial expression.
\end{CoreProblem}
\begin{CoreProblem}{C}{Constrained counting}
{\scriptsize\textit{Name and justify the model; exact fractions and unsimplified products are acceptable.}\par}
How many 5-card hands from a standard 52-card deck contain exactly
    two hearts? Give an unsimplified product of binomial coefficients.
\end{CoreProblem}
\columnbreak
\begin{CoreProblem}{D}{Probability models}
{\scriptsize\textit{Name and justify the model; exact fractions and unsimplified products are acceptable.}\par}
A bag contains 4 red and 6 blue tokens. Two tokens are drawn without
    replacement. Find the probability that the first is red and the second is
    blue. Show the conditional product.
\end{CoreProblem}
\begin{CoreProblem}{E}{Modeling and interpretation}
{\scriptsize\textit{Name and justify the model; exact fractions and unsimplified products are acceptable.}\par}
A condition occurs in $1\%$ of a population. A test is positive for
    $95\%$ of people who have the condition and for $8\%$ of people who do not.
    Find $P(\text{positive})$ and then
    $P(\text{condition}\mid\text{positive})$. Give a one-sentence base-rate
    interpretation.
\end{CoreProblem}
\end{multicols}
\CoreWorkArea
\end{document}
